\section{Comparison of invariant virtual spaces}%
\label{sec:peps_g_injective_virtual_comparison}

\begin{definition}[Virtual comparison]%
    \label{def:peps_g_injective_virtual_comparison}
    \lean{TNLean.PEPS.IsGInjective.virtualEquiv}
    \leanok
    \uses{def:peps_g_injective,thm:peps_g_injective_invariants_range}
    Let $G$ and $H$ be finite groups acting linearly on complex
    vector spaces $\mathcal V$ and $\mathcal W$, respectively.
    Let $T:\mathcal V\to\mathcal P$ be $G$-injective and let
    $S:\mathcal W\to\mathcal P$ be $H$-injective, with
    $\operatorname{im}T=\operatorname{im}S$.
    The virtual comparison is the linear isomorphism
    \begin{align}
        \Phi=(S|_{\mathcal W^H})^{-1}\circ T|_{\mathcal V^G}
        :\mathcal V^G\longrightarrow\mathcal W^H.
    \end{align}
    Both restrictions are isomorphisms onto the common physical range.
\end{definition}

\begin{theorem}[Uniqueness of virtual comparison]%
    \label{thm:peps_g_injective_virtual_comparison}
    \lean{TNLean.PEPS.IsGInjective.virtualEquiv_apply,
        TNLean.PEPS.IsGInjective.virtualEquiv_unique,
        TNLean.PEPS.IsGInjective.virtualEquiv_symm,
        TNLean.PEPS.IsGInjective.existsUnique_virtualEquiv,
        TNLean.PEPS.IsGInjective.finrank_invariants_eq_of_range_eq}
    \leanok
    \uses{def:peps_g_injective_virtual_comparison}
    Under the hypotheses of
    Definition~\ref{def:peps_g_injective_virtual_comparison},
    $\Phi$ is the unique linear map
    $F:\mathcal V^G\to\mathcal W^H$ satisfying $S(Fx)=Tx$
    for every $x\in\mathcal V^G$.
    Reversing $T$ and $S$ gives $\Phi^{-1}$.
    In finite dimension, $\dim\mathcal V^G=\dim\mathcal W^H$.
    This is a local consequence of the invariant-space injectivity in
    \cite[Definition~5.1]{Schuch2010PEPS}.
\end{theorem}

\begin{proof}\leanok
    \uses{thm:peps_g_injective_invariants_range}
    The defining composition gives $S(\Phi x)=Tx$.
    If $S(Fx)=Tx$, then $S(Fx)=S(\Phi x)$, and injectivity of
    $S|_{\mathcal W^H}$ gives $Fx=\Phi x$.
    Applying the same characterization with the two maps reversed
    identifies the inverse. The dimension equality follows from the
    linear isomorphism.
\end{proof}
